\section{The Tier Model}
\label{sec:tiermodel}

We present the Unified Tier Model of Creative Quality: a three-tier hierarchical
structure governing the contribution of prompt components to creative output quality,
identified empirically across authoring and reviewing roles and formalized here as a
generalizable architecture.

The model organizes the components of a creative AI prompt into three tiers by the
functional role each plays in quality production. Tier 1 (Orientation) components are
irreplaceable: removing them causes categorical failure or sub-baseline performance.
Tier 2 (Lens) components are load-bearing: removing either of two identified
components drops quality below the publishable threshold, but their positive effect
depends on the presence of Tier 1. Tier 3 (Specificity) components are amplifying:
each adds marginal quality beyond the Tier 1+2 base, but none is necessary for
threshold crossing, and their aggregate effect is bounded. The tiers are not additive
in a flat sense --- they are ordered in a way that determines which components can
function at all.

\subsection{Tier 1: Orientation (Irreplaceable)}

Tier 1 is defined as a stance toward the person on the other end of the creative
interaction. In authoring, this is the solver who will encounter the puzzle. In
reviewing, this is the reader of the review --- the team that will act on its
recommendations. In both roles, the orientation takes the form of a single directing
question: what does the person on the other end experience? This question precedes and
governs every subsequent creative decision.

The reviewing evidence is provided by the C4 and C5 conditions in the expert panel
study. C4 tested the design philosophy and worldview component in isolation, with no
review lens and no credentials, achieving 26.7 out of 30. C5 tested the review lens
in isolation --- the defect-detection specificity component that defines the reviewer's
analytical focus --- with no philosophy, achieving 21.7. The bare baseline with no
persona at all (C0) scored 23.7. The lens-without-philosophy condition therefore fell
below no-persona: adding domain-specific focus without orientation produced outputs
worse than unstructured generation. The full principles condition (C3) scored 26.7,
identical to C4, suggesting that the eleven principles collectively instantiate the
same orientation that the philosophy states directly. The philosophy is not additive
to the lens; it is the condition under which the lens produces functional outputs
rather than noise.

The authoring evidence is provided by the L-T1 ablation in the trait decomposition
study. The full artist profile scored 41.7. Removing the experience-first trait (T1)
while retaining all five other traits produced a score of 12.3 --- a drop of 29.4
points. No other single-trait removal caused a drop of more than 12.7 points, and
no other removal produced categorical failure. The L-T1 output did not produce a
degraded puzzle; it produced a list of Age of Empires~II terms with no clue layer,
no solver path, and no aha moment. The puzzle ceased to exist as a
solver-experience artifact. The five remaining traits --- elegance drive, visual
thinking, surprise, rigor, systematic reasoning --- produced domain-appropriate format
with no experiential content, because there was no organizing question of what the
solver would experience.

The mechanism is the same in both roles. Orientation establishes what counts as a
successful output. Without it, the system optimizes for an undefined objective:
domain-appropriate format, reviewer-appropriate vocabulary, the surface grammar of
the artifact type. The remaining components then produce outputs that satisfy format
constraints while failing the experience constraint that is the actual quality
criterion. Tier 1 is irreplaceable because it is the criterion itself, not a
contributor to the criterion.

\subsection{Tier 2: Lens (Load-bearing)}

Tier 2 is defined as operational constraints that enforce the Tier 1 orientation
across specific decision points. Lens components answer the question: given the
orientation, what are the checkable tests that translate it into creative decisions?
They are functional rather than aspirational --- they do not state a goal but specify
a procedure that must be followed to achieve the goal.

The reviewing evidence identifies three non-redundant principles from the full set of
eleven. Verify the Last Mile requires the reviewer to confirm that the solver can
complete the final extraction step without implicit author knowledge. Blame the Player
requires the reviewer to consider whether a reported solver failure reflects design
ambiguity before attributing it to solver error. Snyder's Computer Test requires the
reviewer to assess whether the puzzle mechanism could be completed by a mechanical
process with no human judgment, as a test of whether genuine insight has been
required. These three principles are non-redundant in the sense that each catches
defect classes that the others miss and that the full profile's worldview alone does
not surface. The other eight principles --- including Solving Equals Proving
Understanding, No Over-Scaffolding, and the Riven Standard --- are subsumed by the
philosophy: a reviewer who has genuinely internalized the worldview will apply these
principles implicitly without being instructed to do so.

The load-bearing designation follows from the C7 result. When all eleven principles
were added to the full profile (C6 = 25.3), quality dropped to 22.3. Discipline
without orientation produces noise, not signal: principles applied without a
governing worldview generate contradictory constraints and conservative hedging
rather than directed assessment. When the three non-redundant principles are added to
the full profile, quality reaches the best observed condition. The principles are
load-bearing only in the presence of the Tier 1 orientation; absent it, they are
harmful.

The authoring evidence identifies two load-bearing traits. Elegance drive (T4) is
the disposition to remove any element of the puzzle mechanism that does not contribute
to the aha moment: it operationalizes the Tier 1 experience-first orientation into a
cutting test applied at each design decision. Rigor and verification (T5) is the
disposition to test every step of the solver path for unintended solutions, ambiguous
extractions, and knowledge-assumption failures: it operationalizes the orientation
into a verification procedure. Removing T4 dropped quality to 30.0, below the 35-point
publishable threshold (The rubric has seven dimensions: Clarity $\times 1$, Solvability
$\times 1$, Elegance $\times 2$, Reading Reward $\times 2$, Fun $\times 1$, Confirmation
$\times 1$, Riven Standard $\times 1$, maximum 45 points; Elegance and Reading Reward are
double-weighted as the empirically strongest tier discriminators; the 33/45 threshold =
73\% of maximum, matching the 22/30 = 73\% threshold in Paper~\ref{paper1}, calibrated
against quality levels at which puzzles in MIT Mystery Hunt and Microsoft Puzzlehunt have
shipped.). Removing T5 dropped quality to 29.0. Both removals kept the
output above the catastrophic failure level of L-T1, but produced puzzles that were
experiential in aspiration but mechanically flawed or unnecessarily complex ---
orientation without discipline.

The T1+T4+T5 reconstruction condition (R3) is the key evidence for load-bearing
status. Three traits from six produced 35.7 --- the first threshold crossing in any
reconstruction condition. The other three traits (visual thinking, surprise,
systematic reasoning) are together responsible for the gap between 35.7 and 41.7,
confirming their Tier 3 amplifying role. Tier 2 is the minimum necessary complement
to Tier 1: orientation states the goal, discipline makes the goal actionable.

\subsection{Tier 3: Specificity (Amplifying)}

Tier 3 is defined as domain-specific knowledge, authority signals, and stylistic
particulars that enrich output quality beyond the Tier 1+2 base. Specificity
components answer the question: given orientation and discipline, what domain content
produces the best possible output in this particular artifact type?

The reviewing evidence identifies two Tier 3 components. The review lens (defect-
detection specificity) provides a structured analytical framework for identifying
classes of puzzle defect --- extraction ambiguity, clue-answer mismatch, solver-path
incompleteness --- that the orientation-only reviewer would identify as problems but
might not name or categorize consistently. The C4 versus C6 comparison shows that
adding the lens to the philosophy produces a 26.7 versus 25.3 comparison, a modest
gain. Credentials and authority signals calibrate scoring conservatism: a reviewer
whose profile specifies deep experience with high-difficulty puzzles applies scoring
rubrics at a different point on the scale than an uncalibrated reviewer. The lens
caught the specific defect in the test puzzle --- a bracket label error --- that
conditions C0 through C5 all missed; credentials prevented grade inflation in
conditions where the Tier 1+2 base was present and functional.

The authoring evidence identifies three Tier 3 traits. Surprise (T2) is the
disposition to seek mechanisms that produce an unexpected but inevitable aha moment;
it adds 5.4 points to the Tier 1+2 base. Visual thinking (T3) is the disposition to
reason about the puzzle as a spatial layout that guides the solver's eye; it adds
4.4 points. Systematic reasoning (T6) is the disposition to enumerate the design space
before committing to a mechanism; it adds 1.6 points. Domain world data --- the
corpus of Age of Empires~II civilizations, units, technologies, and map names that
constitute the puzzle's thematic content --- is also Tier 3: it provides the raw
material for clue construction but does not determine the mechanism quality or the
solver experience quality.

The artist versus puzzle designer finding is the most direct Tier 3 evidence. The
puzzle designer label activates Tier 3 content: puzzle conventions, format
expectations, the genre grammar of crossword-adjacent mechanisms. The artist label
activates Tier 1 content: the orientation toward what the viewer or audience will
experience. The 15.6-point gap between these two single-word labels measures the cost
of starting at the wrong tier --- domain specificity without orientation.

Domain expertise matters, but it matters last and least. It is necessary for ceiling
performance and it is substitutable: a Tier 1+2 prompt with incorrect domain content
still crosses the quality threshold, while a Tier 3 prompt without Tier 1 does not.

\subsection{The Symmetric Structure}

Table~\ref{tab:tier-symmetric} presents the tier structure symmetrically across
authoring and reviewing roles. The symmetry is structurally informative: it shows that
the tier model is not an artifact of the authoring domain, the reviewing domain, or
the specific puzzle-hunt testbed, but a property of the creative quality architecture
that governs both roles.

\begin{table}[h]
\centering
\caption{Tier model symmetric across authoring and reviewing roles.}
\label{tab:tier-symmetric}
\begin{tabular}{p{2.2cm} p{5.2cm} p{5.2cm}}
\toprule
\textbf{Tier} & \textbf{Reviewing} & \textbf{Authoring} \\
\midrule
\textbf{T1: Orientation} (Irreplaceable)
& Design philosophy and worldview. C5 (lens only) $= 21.7$, below C0 baseline $= 23.7$. C4 (philosophy only) $= 26.7$.
& Experience-first stance (T1). L-T1 $= 12.3/45$; catastrophic failure; puzzle ceased to function as solver artifact. \\
\addlinespace
\textbf{T2: Lens} (Load-bearing)
& 3 non-redundant principles: Verify the Last Mile, Blame the Player, Snyder's Computer Test. C7 (all 11 principles) $= 22.3 <$ C6 $= 25.3$.
& Elegance drive (T4) + Rigor/verification (T5). R3 (T1+T4+T5) $= 35.7$, first threshold crossing. \\
\addlinespace
\textbf{T3: Specificity} (Amplifying)
& Review lens + credentials. Lens caught bracket error missed by C0--C5; credentials calibrate conservatism.
& Surprise (T2), visual thinking (T3), systematic reasoning (T6) + domain world data. Net amplification: $+6.0$ above R3. \\
\bottomrule
\end{tabular}
\end{table}

The connections to Csikszentmihalyi's creator/domain/field distinction and Dreyfus's
expertise stages are offered as structural analogues discovered after derivation, not
causal derivations; the tier model emerges from ablation data, not from these frameworks.

The symmetry follows from a general principle about creative quality. Quality is
defined from the receiver's perspective: the solver who encounters the puzzle, or the
team that reads the review and acts on it. The definition of quality is therefore a
Tier 1 claim, prior to any domain claim. Quality is enforced through operational
constraints that apply the receiver-first definition at decision points: these are
Tier 2. Quality is enriched by domain knowledge that populates the artifact with
specific, appropriate content: this is Tier 3. The ordering holds regardless of role
because the receiver-first definition is prior to the domain in both cases.

The asymmetry between the tier contents --- philosophy versus experience-first
orientation, principles versus elegance-and-rigor traits --- reflects the different
decision structures of the two roles rather than a difference in the tier model itself.
Reviewers make categorical judgments about completed artifacts; their Tier 2 components
are classification principles (does this defect class apply?). Authors make
generative decisions at each design step; their Tier 2 components are process
constraints (should this element be cut? is this step verified?). The Tier 2 function
--- operationalizing orientation into checkable tests --- is the same in both roles;
the specific tests differ because the creative decisions differ.

The OLE formula formalizes this architecture as a prescriptive ordering for AI
creative prompt design. Section~\ref{sec:ole} develops the formula and its
generalization beyond the puzzle-hunt testbed.
